% Teil II: durchgerechneter Weg, Teile a--g.

\subsection*{Gegeben und gesucht}
Kugel $R$, $\rho_{\mathrm K}$ an Feder $D$, in Wasser ($\rho_{\mathrm W}$,
Viskosität $\eta$). Gesucht: (a) $x_R$, (b) $\alpha$, (c) DGL, $\omega_0$,
$\beta$, (d) $R$ für periodische Schwingung, (e) Energieverlust pro Oszillation,
(f) DGL des getriebenen Systems, (g) $x_0/z_0$ und Phase bei $\omega_0$.

\subsection*{Ansatz}
(a) Kräftegleichgewicht in Ruhe (mit Auftrieb). (b) Stoke'sche Kugelreibung.
(c) DGL um die Ruhelage, $\omega_0^2=D/m$, $\beta=\alpha/2m$. (d) unterkritisch
$\beta<\omega_0$ mit $m=\frac43\pi R^3\rho_{\mathrm K}$. (e) $E(t)=E_0e^{-2\beta t}$.
(f)--(g) getriebener Oszillator mit $z(t)$; Ausschwingverhalten aus der
Übertragungsfunktion bei $\omega_0$.

\subsection*{Rechnung}
\textbf{a) Statische Dehnung.} Gewicht $-$ Auftrieb $=$ Federkraft:
\begin{equation}
  D\,x_R=\rho_{\mathrm K}Vg-\rho_{\mathrm W}Vg
  =(\rho_{\mathrm K}-\rho_{\mathrm W})\,\tfrac43\pi R^3\,g
  \;\Rightarrow\;\boxed{\,x_R=\frac{(\rho_{\mathrm K}-\rho_{\mathrm W})\tfrac43\pi R^3 g}{D}\;}.
\end{equation}

\textbf{b) Reibungskoeffizient (laminar).} Stokessche Kugelformel:
\begin{equation}
  \boxed{\;\alpha=6\pi\eta R\;},\qquad F=-\alpha v\;.
\end{equation}

\textbf{c) Bewegungsgleichung und Kenngrössen.} Um die Ruhelage (Auftrieb und
$G$ heben sich auf):
\begin{equation}
  m\,\ddot x=-D\,x-\alpha\,\dot x
  \;\Rightarrow\;
  \boxed{\;\ddot x+2\beta\,\dot x+\omega_0^2 x=0\;},
  \qquad
  \boxed{\;\omega_0=\sqrt{\frac{D}{m}}\;},\quad
  \boxed{\;\beta=\frac{\alpha}{2m}=\frac{3\pi\eta R}{m}\;}.
\end{equation}

\textbf{d) Periodizitätsbedingung.} Unterkritisch $\beta<\omega_0$. Mit
$m=\frac43\pi R^3\rho_{\mathrm K}$:
\begin{equation}
  \frac{3\pi\eta R}{m}<\sqrt{\frac{D}{m}}
  \;\Longleftrightarrow\;
  9\pi^2\eta^2 R^2<D\,m=\frac43\pi\rho_{\mathrm K}D\,R^3
  \;\Rightarrow\;
  \boxed{\;R>\frac{27\,\pi\,\eta^2}{4\,\rho_{\mathrm K}D}\;}.
\end{equation}
Zu kleine Kugeln ($\beta>\omega_0$) kriechen aperiodisch; ab einem
Mindestradius $R_{\min}=27\pi\eta^2/(4\rho_{\mathrm K}D)$ schwingt das Pendel.

\textbf{e) Energieverlust in der ersten Oszillation (schwache Dämpfung).}
$E(t)=E_0e^{-2\beta t}$; nach einer Periode $T_0\approx2\pi/\omega_0$:
\begin{equation}
  \frac{\Delta E}{E_0}=1-e^{-2\beta T_0}
  \;\approx\;\boxed{\;2\beta T_0\;}
  \qquad(\text{klein für } \beta T_0\ll1)\;.
\end{equation}

\textbf{f) Getriebenes System.} Aufhängung bewegt sich um $z(t)$; die
Federkraft ist $-D(x-z)$, die Reibung $- \alpha\dot x$ (Wasser ruht):
\begin{equation}
  m\,\ddot x=-D(x-z)-\alpha\,\dot x
  \;\Longrightarrow\;
  \boxed{\;\ddot x+2\beta\,\dot x+\omega_0^2 x=\omega_0^2\, z(t)\;}.
\end{equation}

\textbf{g) Resonanz: Amplituden- und Phasenverhältnis.} Mit $z(t)=z_0e^{i(\omega_0 t+\varphi)}$
und $x(t)=x_0e^{i\omega_0t}$ ist die Übertragungsfunktion bei $\omega=\omega_0$:
\begin{equation}
  H(\omega_0)=\frac{x_0}{z_0}=
  \frac{\omega_0^2}{\omega_0^2-\omega_0^2+2i\beta\omega_0}
  =\frac{\omega_0}{2\beta}\,e^{-i\pi/2}\;,
\end{equation}
also
\begin{equation}
  \boxed{\;\frac{x_0}{z_0}=\frac{\omega_0}{2\beta}=Q\;},\qquad
  \boxed{\;\text{Phasendifferenz }\delta=\frac{\pi}{2}\;}
  \quad(\text{Auslenkung hinkt der Anregung nach}).
\end{equation}
Mit $\beta=\frac{9\pi\eta}{4\rho_{\mathrm K}R^2}\propto R^{-2}$ wird für
\emph{grösseres} $R$ die Dämpfung kleiner und damit $Q$ grösser; für dieselbe
Auslenkung $x_0$ braucht man dann eine \boxed{\;\emph{kleinere}\;} Amplitude $z_0$.

\begin{ergebnis}
$x_R=\frac{(\rho_{\mathrm K}-\rho_{\mathrm W})Vg}{D}$, $\alpha=6\pi\eta R$,
$\omega_0=\sqrt{D/m}$, $\beta=3\pi\eta R/m$,
$R>27\pi\eta^2/(4\rho_{\mathrm K}D)$,
$\Delta E/E_0\approx2\beta T_0$,
$\ddot x+2\beta\dot x+\omega_0^2x=\omega_0^2 z$,
$x_0/z_0=\omega_0/2\beta$, $\delta=\pi/2$; grösseres $R\Rightarrow$ kleineres $z_0$.
\end{ergebnis}

\begin{kontrolle}
\textbf{Einheiten:} $x_R$: $\dfrac{\mathrm{kg/m^3\cdot m^3\cdot m/s^2}}{N/m}=\dfrac{\mathrm{kg\,m\,s^{-2}}}{N/m}=\mathrm{m}$. \checkmark
$\alpha=6\pi\eta R$: $\mathrm{Pa\,s\cdot m}=kg/s$ --- es ist $[F]=[\alpha][v]$, also
$[\alpha]=[F]/[v]=\tfrac{kg\,m/s^2}{m/s}=kg/s$. \checkmark
$R_{\min}$: $\dfrac{\eta^2}{\rho D}$:
$\dfrac{(kg\,m^{-1}s^{-1})^2}{kg/m^3\cdot kg\,s^{-2}}=\mathrm m$. \checkmark

\textbf{Grenzfall:} $\eta\to0\Rightarrow R_{\min}\to0$: beliebige Kugel schwingt
ungedämpft --- plausibel. $Q=\omega_0/2\beta$: grosse $Q$ $\Leftrightarrow$
schwache Dämpfung, grosse Resonanzüberhöhung. \checkmark
\end{kontrolle}

\begin{merksatz}
Die Periodizitätsbedingung $\beta<\omega_0$ ist hier eine \emph{Minimal-}Bedingung
an $R$, weil die Kugelmasse mit $R^3$ wächst: mehr Masse überwiegt die
lineare Stokes-Reibung. Bei Resonanz ist die Auslenkung um $\pi/2$ verzögert
und um $Q=\omega_0/2\beta$ überhöht.
\end{merksatz}
